%=========================================================
% TRABALHO DE CONCLUSÃO DE CURSO
% Faculdade Nova Roma — Ciência da Computação
%
% Documento mestre. As seções estão em sections/ e as figuras/tabelas
% são geradas pelo pipeline em ../data/analytics.
% Compilar:  pdflatex main.tex   (duas vezes, para o sumário/refs)
%=========================================================
\documentclass[a4paper,12pt]{extarticle}

\usepackage[brazil]{babel}
\usepackage[utf8]{inputenc}
\usepackage[T1]{fontenc}

% Fonte sem serifa (padrão visual do documento institucional)
\usepackage[scaled]{helvet}
\renewcommand{\familydefault}{\sfdefault}

\usepackage{graphicx}
\usepackage{tikz}
\usetikzlibrary{arrows.meta,positioning,shapes.geometric,calc,fit,backgrounds}
\usepackage{eso-pic}
\usepackage[
    a4paper,
    top=2.5cm,
    bottom=2.5cm,
    left=3cm,
    right=3cm
  ]{geometry}
\usepackage{booktabs}
\usepackage{float}
\usepackage{setspace}
\usepackage{indentfirst}
\usepackage{titlesec}
\usepackage{caption}
\captionsetup{font=small,labelfont=bf}
\usepackage[hidelinks]{hyperref}

% Localiza as figuras geradas pelo pipeline de análise.
\graphicspath{{figures/}}

%--------------------------------------------------------
% Dados do trabalho

%--------------------------------------------------------
% Dados do trabalho
%--------------------------------------------------------
\newcommand{\tccautor}{Maria Júlia Oliveira}
\newcommand{\tcctitulo}{Exposição Gradual a Estímulos Fóbicos em Realidade Virtual: desenvolvimento e validação de um protótipo web}
\newcommand{\tccorientador}{Prof.\ MSc.\ Henning Summer}
\newcommand{\tcclocal}{Recife}
\newcommand{\tccano}{2026}

%--------------------------------------------------------
% Formatação de títulos (numerados, em negrito; seções em caixa alta)
%--------------------------------------------------------
\titleformat{\section}{\normalfont\bfseries}{\thesection.}{0.5em}{\MakeUppercase}
\titleformat{\subsection}{\normalfont\bfseries}{\thesubsection}{0.5em}{}
\titleformat{\subsubsection}{\normalfont\bfseries\itshape}{\thesubsubsection}{0.5em}{}
\titlespacing*{\section}{0pt}{1.4\baselineskip}{0.7\baselineskip}
\titlespacing*{\subsection}{0pt}{1.1\baselineskip}{0.5\baselineskip}
\titlespacing*{\subsubsection}{0pt}{0.9\baselineskip}{0.4\baselineskip}

% Espaçamento entrelinhas (1,5) e recuo de parágrafo
\onehalfspacing
\setlength{\parindent}{1.5cm}

%--------------------------------------------------------
% TIMBRE — aplicado a TODAS as páginas (papel timbrado institucional)
%--------------------------------------------------------
\AddToShipoutPictureBG{%
  \begin{tikzpicture}[remember picture,overlay]
    % Marca d'água lateral esquerda
    \node[anchor=north west, inner sep=0, xshift=-10mm, yshift=10mm]
      at (current page.north west) {\includegraphics{timbre/background-nr.png}};
    % Logo superior direita
    \node[anchor=north east, inner sep=1.3cm, yshift=5mm]
      at (current page.north east) {\includegraphics{timbre/logo-nr.png}};
  \end{tikzpicture}%
}

\begin{document}
\pagenumbering{gobble}

%=================== ELEMENTOS PRÉ-TEXTUAIS ===================
%=========================== CAPA ===========================
\begingroup
\singlespacing
\begin{center}
\vspace*{1.2cm}

{\bfseries FACULDADE NOVA ROMA}\\[2pt]
{\bfseries CURSO DE CIÊNCIA DA COMPUTAÇÃO}

\vfill

{\bfseries \tccautor}

\vfill

{\large\bfseries \tcctitulo}

\vfill
\vfill

{\bfseries \MakeUppercase{\tcclocal}}\\[2pt]
{\bfseries \tccano}

\vspace*{0.6cm}
\end{center}
\endgroup
\newpage

%======================= FOLHA DE ROSTO =======================
\begingroup
\singlespacing
\begin{center}
\vspace*{1.2cm}

{\bfseries \tccautor}

\vfill

{\large\bfseries \tcctitulo}
\end{center}

\vspace{1.2cm}

% Natureza do trabalho — bloco recuado à direita (padrão ABNT)
\noindent\hspace*{7cm}\begin{minipage}{7.5cm}
\singlespacing\small
Trabalho de Conclusão de Curso apresentado como requisito para aprovação na
disciplina de Trabalho de Conclusão de Curso~II do curso de Ciência da
Computação da Faculdade Nova Roma (Caruaru/Recife).

\vspace{0.8cm}
Orientador: \tccorientador
\end{minipage}

\vfill

\begin{center}
{\bfseries \MakeUppercase{\tcclocal}}\\[2pt]
{\bfseries \tccano}
\end{center}
\vspace*{0.6cm}
\endgroup
\newpage

%===================== FOLHA DE APROVAÇÃO =====================
\begingroup
\singlespacing
\begin{center}
\vspace*{1.2cm}

{\bfseries \tccautor}

\vfill

{\large\bfseries \tcctitulo}
\end{center}

\vspace{1.0cm}

\hspace*{7cm}\begin{minipage}{7.5cm}
\singlespacing\small
Trabalho de Conclusão de Curso apresentado como requisito para aprovação na
disciplina de Trabalho de Conclusão de Curso~II do curso de Ciência da
Computação da Faculdade Nova Roma (Caruaru/Recife).
\end{minipage}

\vspace{1.4cm}
\noindent Aprovado em \rule{1.2cm}{0.4pt}\,/\,\rule{1.2cm}{0.4pt}\,/\,\rule{1.8cm}{0.4pt}.

\vspace{1.6cm}
\begin{center}
{\bfseries BANCA EXAMINADORA}
\end{center}

\vspace{1.4cm}
\begin{center}
\rule{10cm}{0.4pt}\\
\tccorientador\ (Orientador)\\
Faculdade Nova Roma

\vspace{1.4cm}
\rule{10cm}{0.4pt}\\
Prof.\ \rule{5cm}{0.4pt} (Examinador(a))\\
Faculdade Nova Roma

\vspace{1.4cm}
\rule{10cm}{0.4pt}\\
Prof.\ \rule{5cm}{0.4pt} (Examinador(a))\\
Faculdade Nova Roma
\end{center}

\vfill
\endgroup
\newpage

%=========================== RESUMO ===========================
\begingroup
\begin{center}{\bfseries RESUMO}\end{center}
\vspace{0.3cm}

% ESCREVER POR ÚLTIMO, depois dos resultados. O que está aqui já
% segue a fórmula de cinco movimentos: contexto, impasse, o que foi
% feito, como foi avaliado, o que se achou. Os dois últimos só podem
% ser fechados após a coleta.

\noindent Fobias específicas a animais impõem limitações cotidianas a
parcela expressiva da população, e a superação desse tipo de medo depende
do contato progressivo com o estímulo temido --- que é justamente aquilo
que a pessoa evita. Ambientes virtuais oferecem uma via para contornar
esse impasse, permitindo graduar o estímulo enquanto o controle permanece
com quem é exposto. Este trabalho apresenta o desenvolvimento e a
validação de um protótipo de Realidade Virtual executável em navegador,
sem instalação, que simula cinco animais associados a fobias --- cobra,
rato, barata, aranha e sapo --- cada um no ambiente em que seria
encontrado, com a exposição organizada em cinco níveis progressivos nos
quais variam a distância, a quantidade de exemplares, o movimento e a
possibilidade de interação direta. A aplicação registra automaticamente
medidas de comportamento durante o uso, entre elas a menor distância
alcançada, a proporção do tempo com o estímulo no campo de visão e o
nível máximo atingido, confrontadas com o grau de medo declarado
previamente pelo participante. [COMPLETAR: a validação foi conduzida
com N participantes dos cursos de X e Y.] [COMPLETAR: o que os dados
indicaram.] Ressalta-se que o sistema constitui protótipo acadêmico, não
se destinando a tratamento ou diagnóstico.

\vspace{0.4cm}
\noindent\textbf{Palavras-chave:} Realidade Virtual; Fobias específicas;
Exposição gradual; WebXR; Usabilidade.

\vspace{1.0cm}
\begin{center}{\bfseries ABSTRACT}\end{center}
\vspace{0.3cm}

% Traduzir depois que o português estiver fechado.
\noindent Specific animal phobias impose everyday limitations on a
significant portion of the population, and overcoming such fears depends
on progressive contact with the feared stimulus --- precisely what the
person avoids. Virtual environments offer a way around this impasse,
allowing the stimulus to be graded while control remains with the person
being exposed. This work presents the development and validation of a
browser-based Virtual Reality prototype, requiring no installation, that
simulates five phobia-associated animals --- snake, rat, cockroach,
spider and frog --- each in the environment where it would be
encountered, with exposure organized into five progressive levels varying
distance, number of specimens, movement and the possibility of direct
interaction. The application automatically records behavioural measures
during use, including the closest distance reached, the proportion of
time with the stimulus in the field of view and the highest level
attained, which are compared against the degree of fear previously
declared by the participant. [TO COMPLETE.] The system is an academic
prototype and is not intended for treatment or diagnosis.

\vspace{0.4cm}
\noindent\textbf{Keywords:} Virtual Reality; Specific phobias; Gradual
exposure; WebXR; Usability.
\endgroup
\newpage


\newpage
\tableofcontents
\newpage

%=================== ELEMENTOS TEXTUAIS ===================
\pagenumbering{arabic}
%========================= INTRODUÇÃO =========================
\section{Introdução}

% RASCUNHO. A estrutura e o fio do argumento estão prontos; falta
% inserir as citações, que só podem entrar depois da revisão.
% Onde está [CITAR], entra a referência do artigo correspondente.

Fobias específicas estão entre as condições de maior prevalência
relacionadas à ansiedade, e as associadas a animais figuram entre as mais
comuns delas [CITAR]. Cobras, ratos, baratas, aranhas e sapos aparecem de
forma recorrente nessa categoria --- animais que, no contexto urbano,
raramente oferecem risco objetivo, mas cuja simples possibilidade de
encontro basta para alterar a rotina de quem os teme.

A limitação que daí decorre é menos visível que a reação em si. Ela se
manifesta em evitação: deixar de entrar num cômodo, de visitar um lugar,
de exercer uma atividade. [CITAR]

\subsection*{O impasse}

O procedimento reconhecido para a superação desse tipo de medo envolve o
contato progressivo e controlado com o estímulo temido [CITAR]. E é
justamente aí que se instala um impasse: \emph{o contato é precisamente
aquilo que a pessoa evita}. O procedimento depende daquilo que a
condição impede.

Ambientes virtuais oferecem uma via para contornar esse impasse. Neles, o
estímulo pode ser graduado com precisão --- aproximado, multiplicado,
posto em movimento --- enquanto o participante mantém a certeza de que
nada ali é real e de que pode interromper a qualquer instante. O controle
permanece com quem é exposto, o que é condição do procedimento e não um
detalhe de conforto.

\subsection*{A lacuna}

% COMPLETAR com o que a revisão efetivamente mostrar. A redação abaixo
% é a hipótese de partida e precisa ser confirmada ou corrigida.
A literatura que investiga exposição a estímulos fóbicos em Realidade
Virtual concentra-se, em larga medida, em estudos conduzidos em
laboratório, com equipamento dedicado e acompanhamento presencial
[CITAR]. São condições que asseguram controle experimental, mas que
restringem o alcance: dependem de um local específico, de um aparelho
específico e de uma pessoa presente para operá-lo.

A consolidação do WebXR --- especificação que permite executar conteúdo
tridimensional imersivo diretamente no navegador --- altera esse quadro,
ao viabilizar que a mesma aplicação funcione em computador, celular e
óculos sem instalação. Discute-se menos, contudo, o que essa mudança
acarreta: como as pessoas percebem e interagem com estímulos fóbicos
quando o ambiente virtual é alcançado por um endereço, e não por um
laboratório.

\subsection*{Questão de pesquisa e objetivo}

Deste trabalho, portanto, a questão:

\begin{quote}
Como as pessoas percebem e interagem com estímulos associados a fobias de
animais em ambientes virtuais organizados em níveis progressivos de
aproximação?
\end{quote}

Para investigá-la, foi desenvolvido um protótipo que simula cinco animais
associados a fobias, cada um no ambiente em que seria encontrado, com a
exposição organizada em cinco níveis nos quais variam a distância, a
quantidade de exemplares, o movimento e a possibilidade de interação. A
aplicação registra, durante o uso, medidas de comportamento que são
confrontadas com o grau de medo que o participante declara previamente.

\subsection*{Delimitação de escopo}

Cabe esclarecer, desde a introdução, que o sistema desenvolvido é um
\textbf{protótipo acadêmico}. Ele não constitui tratamento psicológico ou
médico, não realiza diagnóstico e não se propõe substituto de
acompanhamento profissional. A avaliação conduzida restringe-se a
usabilidade, imersão, interação e percepção do estímulo.

Essa delimitação não é ressalva protocolar: ela determina o método. É por
ela que não se aplica instrumento clínico de rastreio, que não há
critério de inclusão por diagnóstico e que o grau de medo é tratado como
autorrelato de percepção, e não como medida clínica.

\subsection*{Organização do texto}

A Seção~2 apresenta os objetivos. A Seção~3 reúne o referencial teórico
construído a partir da revisão da literatura. A Seção~4 descreve a
metodologia, incluindo o protocolo de busca, o perfil dos participantes e
as variáveis coletadas. A Seção~5 descreve o artefato desenvolvido e as
decisões de projeto que o determinaram. A Seção~6 apresenta e discute os
resultados. A Seção~7 conclui, separando o que os dados sustentam do que
permanece como conjectura.

%========================== OBJETIVOS ==========================
\section{Objetivos}

\subsection{Objetivo geral}

% RASCUNHO -- revisar a redação
Desenvolver e validar um protótipo de Realidade Virtual executável em
navegador, voltado à apresentação gradual de estímulos associados a fobias
específicas de animais, avaliando a sua usabilidade, a imersão percebida e a
forma como os participantes interagem com o estímulo ao longo de níveis
progressivos de aproximação.

\subsection{Objetivos específicos}

\begin{itemize}
  \item revisar a literatura sobre Realidade Virtual aplicada à exposição a
  estímulos fóbicos, por meio de busca sistematizada em bases científicas;

  \item projetar e implementar os ambientes virtuais correspondentes a cinco
  animais associados a fobias, bem como a progressão em cinco níveis de
  aproximação;

  \item desenvolver a camada de persistência e autenticação responsável pelo
  cadastro dos participantes e pelo registro das métricas de uso;

  \item coletar e analisar métricas comportamentais de interação e respostas
  de questionários de percepção aplicados a estudantes de Ciência da
  Computação e de Psicologia;

  \item identificar pontos fortes, limitações e oportunidades de melhoria da
  solução, fundamentando recomendações para trabalhos futuros.
\end{itemize}

%===================== REFERENCIAL TEÓRICO =====================
\section{Referencial teórico}

% ESTA É A SEÇÃO QUE O ORIENTADOR MAIS COBROU.
% Ela se escreve a partir dos (até) 10 artigos selecionados na busca.
% Para CADA artigo, anotar: o que mediram, com quantas pessoas, que
% resultado deu. Depois agrupar por tema nas subseções abaixo.

\subsection{Fobias específicas a animais}

% O que são, prevalência, critérios. Base conceitual.
[ESCREVER]

\subsection{Exposição gradual}

% O princípio: contato progressivo, controlado, com o estímulo temido.
% Por que a graduação importa e por que o susto é contraproducente.
[ESCREVER]

\subsection{Realidade Virtual aplicada à exposição}

% O núcleo da revisão. Aqui entram os estudos encontrados: o que já foi
% feito, com que equipamento, com que resultado.
[ESCREVER]

\subsection{Realidade Virtual no navegador}

% WebXR, A-Frame, execução sem instalação. É aqui que se argumenta a
% contribuição deste trabalho frente ao que a literatura já cobriu.
[ESCREVER]

\subsection{Síntese dos trabalhos relacionados}

% Fechar com uma tabela comparativa: trabalho | estímulo | equipamento |
% n | o que mediram. É o que a banca mais olha.
[ESCREVER]

%========================= METODOLOGIA =========================
\section{Metodologia}

% RASCUNHO gerado a partir do código e da documentação do protótipo.
% Revisar, ajustar os números que ainda dependem da coleta e inserir
% as citações dos artigos selecionados na revisão.

Este trabalho caracteriza-se como uma pesquisa aplicada, de natureza
exploratória e abordagem predominantemente quantitativa, organizada em quatro
etapas: revisão da literatura, desenvolvimento do artefato, coleta de dados com
participantes e análise dos resultados.

\subsection{Revisão da literatura}

A fundamentação teórica foi construída a partir de uma busca sistematizada nas
bases Scopus e SpringerLink, restrita a publicações dos últimos cinco anos. A
expressão de busca combinou quatro blocos de termos --- tecnologia, condição,
estímulo e procedimento --- conforme apresentado a seguir:

\begin{quote}\small\singlespacing
\texttt{( "virtual reality" OR "immersive environment" OR WebXR )}\\
\texttt{AND ( phobia OR "specific phobia" OR fear OR anxiety )}\\
\texttt{AND ( spider OR snake OR cockroach OR insect OR animal}\\
\texttt{\phantom{AND (} OR arachnophobia OR ophidiophobia OR entomophobia )}\\
\texttt{AND ( exposure OR "exposure therapy" OR desensitization OR habituation )}
\end{quote}

% TODO: preencher depois de rodar a busca
Os resultados foram filtrados por leitura de título e resumo, chegando-se a um
conjunto final de \textbf{[N]} artigos, que fundamentam as decisões de projeto
descritas na Seção~\ref{sec:desenvolvimento} e a interpretação dos resultados.

\subsection{Participantes}

% TODO: confirmar números após a coleta
A validação foi conduzida com estudantes dos cursos de Ciência da Computação e
de Psicologia da Faculdade Nova Roma. A participação foi voluntária e mediante
assinatura do Termo de Consentimento Livre e Esclarecido (Apêndice~A).

Cabe registrar que \textbf{o sistema não constitui tratamento psicológico ou
médico e não realiza diagnóstico}. A avaliação conduzida restringe-se a
usabilidade, imersão, interação e percepção do estímulo, não se propondo a medir
desfecho clínico de qualquer natureza.

\subsection{Procedimento}

Cada sessão seguiu o mesmo roteiro:

\begin{enumerate}
  \item apresentação do estudo e leitura do termo de consentimento;
  \item cadastro na plataforma, com registro de e-mail, sexo e data de
  nascimento;
  \item questionário inicial de perfil, no qual o participante declara, em
  escala de 0 a 10, o grau de medo que atribui a cada um dos cinco animais
  contemplados;
  \item escolha do estímulo e execução da experiência, percorrendo os níveis de
  exposição na ordem crescente, com liberdade para interromper a qualquer
  momento;
  \item questionário final de percepção.
\end{enumerate}

\subsection{Variáveis coletadas}

Os dados foram registrados automaticamente pela aplicação durante a experiência
e associados a um código de participante (P01, P02, \ldots), de modo que a
análise é conduzida sobre dados anonimizados. O Quadro a seguir relaciona cada
variável comportamental à sua forma de medida.

\begin{table}[htbp]
\centering
\caption{Variáveis comportamentais registradas durante a experiência.}
\label{tab:variaveis}
\begin{tabular}{ll}
\toprule
\textbf{Variável} & \textbf{Como é medida} \\
\midrule
Aproximação          & menor distância alcançada até o estímulo, em metros \\
Afastamento          & número de recuos e distância média mantida \\
Agitação             & variação angular da câmera por segundo \\
Atenção ao estímulo  & proporção do tempo com o estímulo no campo de visão \\
Exploração           & deslocamento acumulado e animais encontrados \\
Interação            & ocorrência de toque no estímulo (Nível 5) \\
Progressão           & nível máximo alcançado e tempo por nível \\
Desconforto          & autorrelato em escala de 0 a 10, por nível \\
\bottomrule
\end{tabular}
\end{table}

\subsection{Tratamento dos dados}

Os registros são gravados em banco de dados relacional e exportados para
análise descritiva. As figuras e tabelas apresentadas na Seção de resultados são
geradas por scripts versionados em \texttt{data/analytics}, a partir dos dados
já anonimizados, o que torna a análise reprodutível sem expor informação
identificável dos participantes.

%======================= DESENVOLVIMENTO =======================
\section{Desenvolvimento}
\label{sec:desenvolvimento}

% Escrita a partir do código. Todos os números desta seção são reais
% e foram medidos. Faltam apenas as capturas de tela, marcadas com
% [FIGURA].

Esta seção descreve o artefato construído e, sobretudo, as decisões de
projeto que o determinaram. A ênfase recai sobre as decisões porque é
nelas que reside a contribuição de engenharia do trabalho: a maior parte
delas não decorre de preferência, mas de restrições concretas ---
perceptivas, de desempenho ou de plataforma --- que só se manifestaram
durante a construção.

\subsection{Tecnologias e organização}

O protótipo é uma aplicação web estática, executável diretamente no
navegador, sem instalação. A camada tridimensional utiliza o A-Frame
(versão 1.7.0), arcabouço declarativo construído sobre o Three.js que
implementa a especificação WebXR. Isso permite que a mesma página
funcione em três modos de uso --- computador, celular e óculos de
Realidade Virtual --- sem base de código separada para cada um.

A opção por HTML, CSS e JavaScript sem arcabouços de interface foi
deliberada, por três razões: reduz a dependência de terceiros, elimina
etapa de compilação (a aplicação é publicada copiando arquivos), e
mantém o código legível como artefato acadêmico --- um leitor que
conheça JavaScript consegue acompanhá-lo sem conhecer um arcabouço
específico.

\begin{table}[htbp]
\centering
\caption{Organização do código-fonte da experiência.}
\label{tab:arquivos}
\begin{tabular}{llr}
\toprule
\textbf{Arquivo} & \textbf{Responsabilidade} & \textbf{Linhas} \\
\midrule
\texttt{dados.js}        & os cinco animais e os cinco níveis            & 321 \\
\texttt{cenarios.js}     & monta os quatro ambientes e os esconderijos   & 597 \\
\texttt{experiencia.js}  & lê a escolha, monta a cena e conduz a sessão  & 1329 \\
\texttt{console.js}      & painel de informação dentro da cena           & 436 \\
\texttt{painel.js}       & painéis de texto desenhados em canvas         & 208 \\
\texttt{controles.js}    & controle de caminhada por toque               & 178 \\
\texttt{animador.js}     & reprodução das animações dos modelos          & 173 \\
\texttt{descoberta.js}   & portas de armário e revelação por aproximação & 136 \\
\midrule
\multicolumn{2}{l}{Total da experiência} & 3.378 \\
\bottomrule
\end{tabular}
\end{table}

O arquivo \texttt{dados.js} ocupa posição particular: ele é a fonte única
de verdade do sistema. Os cinco animais, os cinco níveis, as distâncias,
as quantidades e as escalas estão todos declarados ali, e nenhum outro
arquivo repete esses valores. Alterar a progressão dos níveis é alterar
uma linha --- o que importa num protótipo de pesquisa, em que os
parâmetros do estímulo podem precisar de ajuste entre sessões.

\subsection{Estímulos e ambientes}

Foram contemplados cinco animais frequentemente associados a fobias
específicas. Cada um é apresentado no ambiente em que seria encontrado,
construído com primitivas geométricas do próprio arcabouço.

\begin{table}[htbp]
\centering
\caption{Estímulos, fobias associadas e ambientes de apresentação.}
\label{tab:estimulos}
\begin{tabular}{llll}
\toprule
\textbf{Estímulo} & \textbf{Fobia} & \textbf{Ambiente} & \textbf{Tamanho real} \\
\midrule
Cobra  & Ofidiofobia     & Floresta & 1,40 m \\
Rato   & Musofobia       & Sala     & 0,25 m \\
Barata & Katsaridafobia  & Cozinha  & 0,05 m \\
Aranha & Aracnofobia     & Sala     & 0,08 m \\
Sapo   & Ranidafobia     & Quintal  & 0,10 m \\
\bottomrule
\end{tabular}
\end{table}

% [FIGURA] capturas dos quatro ambientes
A escolha de construir os cenários com primitivas, em vez de importar
ambientes prontos, decorre do peso: os quatro ambientes somados não
acrescentam nenhum arquivo ao projeto, o que mantém o tempo de
carregamento compatível com uma sessão de testes em que cada participante
abre a página do zero.

\subsubsection{Tamanho real e escala de apresentação}

Tamanho real e estímulo perceptível são grandezas distintas. Uma barata
de cinco centímetros, observada a três metros, ocupa menos de cinco
pixels na tela: fisicamente correta e visualmente inútil. A primeira
versão do protótipo apresentava os animais em tamanho real e o efeito
foi exatamente esse --- os estímulos pequenos eram imperceptíveis.

A solução adotada foi declarar separadamente, para cada animal, o tamanho
real e uma escala de apresentação. O sistema apresenta o produto dos
dois, e o trabalho pode informar ambos com transparência.

\begin{table}[htbp]
\centering
\caption{Tamanho real, escala de apresentação e tamanho apresentado.}
\label{tab:escalas}
\begin{tabular}{lrrr}
\toprule
\textbf{Estímulo} & \textbf{Real} & \textbf{Escala} & \textbf{Apresentado} \\
\midrule
Cobra  & 1,40 m & 1,0$\times$ & 1,40 m \\
Rato   & 0,25 m & 1,8$\times$ & 0,45 m \\
Barata & 0,05 m & 4,5$\times$ & 0,23 m \\
Aranha & 0,08 m & 3,6$\times$ & 0,29 m \\
Sapo   & 0,10 m & 3,0$\times$ & 0,30 m \\
\bottomrule
\end{tabular}
\end{table}

Note-se que a escala é inversamente proporcional ao tamanho do animal:
quanto menor o estímulo real, maior o fator necessário para torná-lo
perceptível. A cobra, única de porte comparável ao humano, é apresentada
em tamanho real.

\subsubsection{Ajuste automático dos modelos}

Modelos tridimensionais obtidos de repositórios públicos chegam em
unidades arbitrárias e com a origem em ponto imprevisível. Um mesmo
arquivo pode medir duas ou cento e vinte unidades. O sistema resolve isso
medindo a caixa envolvente do modelo depois do carregamento e calculando
o fator que o leva ao tamanho declarado, além de deslocá-lo
verticalmente para que a base encoste no piso.

Um terceiro ajuste foi necessário e é menos evidente: \emph{a maior
dimensão do modelo nem sempre é aquela a que o tamanho se refere}. O
modelo de cobra empregado inicialmente vinha em pose empinada, de modo
que a sua maior dimensão era a altura; escalar pela maior dimensão
produzia uma cobra de 1,2~m de altura, semelhante a um pilar. Cada animal
passou a declarar, portanto, a qual dimensão o seu tamanho se refere.

Um defeito dessa mesma natureza foi registrado durante o
desenvolvimento e serve de ilustração: a substituição do modelo de cobra
por outro, já em pose estendida, fez a configuração anterior escalar a
dimensão errada, e a verificação automatizada registrou estímulos de 13 a
16 metros de comprimento. O erro foi detectado pela ferramenta, e não a
olho nu.

\subsection{Níveis de exposição}

A progressão é organizada em cinco níveis, nos quais variam quatro
dimensões: a distância até o estímulo, a quantidade de exemplares, a
presença de movimento e a possibilidade de interação direta.

\begin{table}[htbp]
\centering
\caption{Dimensões que variam a cada nível de exposição.}
\label{tab:niveis}
\begin{tabular}{clccc}
\toprule
\textbf{Nível} & \textbf{Denominação} & \textbf{Distância} & \textbf{Movimento} & \textbf{Interação} \\
\midrule
1 & Primeiro contato      & máxima    & não & não \\
2 & Aproximação           & média     & não & não \\
3 & Movimento do estímulo & média     & sim & não \\
4 & Proximidade elevada   & reduzida  & sim & não \\
5 & Interação direta      & mínima    & sim & sim \\
\bottomrule
\end{tabular}
\end{table}

Os níveis 2 e 3 compartilham a mesma distância por definição: entre eles
muda o movimento, não a proximidade. Isolar uma variável por vez é o que
permite, na análise, atribuir eventual mudança de comportamento ao fator
correto.

\subsubsection{Distâncias}

Cada animal possui a sua própria escala de distâncias, porque os tamanhos
variam em uma ordem de grandeza. Cada um declara a distância do Nível~1 e
a do Nível~5; os intermediários são interpolados por uma proporção
declarada no nível.

\begin{table}[htbp]
\centering
\caption{Distância ao estímulo principal, em metros, por nível.}
\label{tab:distancias}
\begin{tabular}{lrrrrr}
\toprule
\textbf{Estímulo} & \textbf{N1} & \textbf{N2} & \textbf{N3} & \textbf{N4} & \textbf{N5} \\
\midrule
Cobra  & 5,50 & 3,61 & 3,61 & 2,14 & 1,30 \\
Rato   & 3,20 & 2,17 & 2,17 & 1,36 & 0,90 \\
Aranha & 2,60 & 1,81 & 1,81 & 1,20 & 0,85 \\
Sapo   & 2,60 & 1,81 & 1,81 & 1,20 & 0,85 \\
Barata & 1,90 & 1,38 & 1,38 & 0,98 & 0,75 \\
\bottomrule
\end{tabular}
\end{table}

\subsubsection{Quantidade}

A quantidade de exemplares cresce de um até quinze, ajustada por um fator
próprio de cada espécie. A necessidade do fator veio da observação: uma
cozinha com vinte e quatro baratas é uma cena reconhecível; uma mata com
quinze cobras de 1,40~m não é --- lê-se como exagero, e não como
intensificação do estímulo.

\begin{table}[htbp]
\centering
\caption{Número de exemplares apresentados por nível.}
\label{tab:quantidades}
\begin{tabular}{lrrrrrr}
\toprule
\textbf{Estímulo} & \textbf{Fator} & \textbf{N1} & \textbf{N2} & \textbf{N3} & \textbf{N4} & \textbf{N5} \\
\midrule
Barata & 1,6 & 1 & 5 & 10 & 16 & 24 \\
Rato   & 1,0 & 1 & 3 & 6  & 10 & 15 \\
Aranha & 1,0 & 1 & 3 & 6  & 10 & 15 \\
Sapo   & 1,0 & 1 & 3 & 6  & 10 & 15 \\
Cobra  & 0,5 & 1 & 2 & 3  & 5  & 8  \\
\bottomrule
\end{tabular}
\end{table}

O Nível~1 ignora o fator: o primeiro contato é sempre com um único
exemplar, para todos os participantes, de modo que seja idêntico e
previsível.

\subsubsection{Reprodutibilidade da disposição}

As posições dos exemplares são sorteadas, mas o sorteio parte de uma
semente derivada do animal e do nível. A consequência é que o mesmo nível
monta sempre a mesma cena. Sem isso, dois participantes veriam
disposições diferentes e a comparação entre eles perderia sentido --- um
poderia ter encontrado três animais a um metro e o outro, nenhum.

\subsubsection{Exclusão deliberada das animações de ataque}

Os modelos empregados trazem animações de ataque e de morte, nomeadas
\texttt{Attack} e \texttt{Death}. Nenhuma delas é utilizada. Um animal
que ataca produz susto, e susto é o oposto de exposição gradual e
controlada: retira do participante o domínio sobre o próprio ritmo, que é
o princípio do procedimento.

A exclusão é uma decisão de projeto fundamentada, e não uma limitação
técnica --- as animações estão nos arquivos e poderiam ser reproduzidas
com uma linha de configuração.

\subsection{Comportamento dos estímulos}

Um grupo em que todos os exemplares se movem de maneira idêntica não
parece um grupo: parece um enfeite. Cada exemplar recebe, portanto, um
comportamento próprio, sorteado com a mesma semente que determina as
posições.

\begin{table}[htbp]
\centering
\caption{Comportamentos atribuídos a cada exemplar.}
\label{tab:comportamentos}
\small
\begin{tabular}{p{2.9cm}p{6.4cm}p{2.8cm}}
\toprule
\textbf{Comportamento} & \textbf{O que faz} & \textbf{Custo} \\
\midrule
Quieto    & apenas respira, com balanço vertical lento & baixo: dispensa o reprodutor de animação \\
Inquieto  & reproduz a animação do arquivo e gira o corpo & médio \\
Andarilho & o mesmo, e ainda caminha devagar entre dois pontos & maior \\
\bottomrule
\end{tabular}
\end{table}

Nos níveis 1 e 2, 65\% dos exemplares permanecem quietos; a partir do
Nível~3, a proporção passa a 25\% quietos, 40\% inquietos e 35\%
andarilhos. Três regras são mantidas em todos os casos:

\begin{enumerate}
  \item \textbf{nenhum exemplar fica completamente imóvel} --- quem não
  reproduz animação recebe ao menos o balanço de respiração;
  \item \textbf{sempre há algum exemplar parado} --- sem contraste, o
  grupo inteiro se torna um borrão de movimento e nada chama atenção;
  \item \textbf{o exemplar principal nunca se desloca} --- ele está na
  distância exata que define o nível, e sair do lugar desfaria a
  progressão que se pretende medir.
\end{enumerate}

\subsection{Exploração e descoberta}

A partir do Nível~2, parte dos exemplares inicia a cena invisível,
alojada em esconderijos distribuídos pelos ambientes. São vinte e um
esconderijos no total --- cinco na sala, nove na cozinha, cinco na
floresta e cinco no quintal --- e a fração de exemplares escondidos é de
35\% nos níveis 2 e 3 e de 45\% nos níveis 4 e 5.

Dois modos de revelação foram implementados:

\begin{table}[htbp]
\centering
\caption{Modos de revelação dos exemplares escondidos.}
\label{tab:revelacao}
\begin{tabular}{p{2.6cm}p{10.0cm}}
\toprule
\textbf{Esconderijo} & \textbf{Como o exemplar aparece} \\
\midrule
Com porta & somente quando o participante \textbf{abre a porta do armário}; a aproximação não basta \\
Sem porta & quando o participante \textbf{chega a menos de 1,6~m} do esconderijo \\
\bottomrule
\end{tabular}
\end{table}

% [FIGURA] cozinha com as portas fechadas e abertas
Dos nove esconderijos da cozinha, cinco são armários com porta articulada
que gira ao ser tocada. Ao aparecer, o exemplar desloca-se para fora do
esconderijo, em direção ao centro do ambiente.

O efeito pretendido é deslocar o controle da exposição. Em lugar de
encarar todos os estímulos de uma só vez, o participante avança no
próprio ritmo e decide quando abrir a próxima porta. O painel de
informação indica quantos exemplares já foram encontrados do total
escondido --- sem esse indicador, quem explora tende a desistir supondo
já ter visto tudo.

\subsection{Interface}
\label{sec:interface}

A interface da aplicação é construída como objeto da cena
tridimensional, e não como camada de HTML sobreposta ao desenho. A
decisão tem motivação técnica direta: ao entrar em modo de Realidade
Virtual, o navegador exibe exclusivamente o conteúdo tridimensional, de
modo que qualquer elemento em HTML desaparece. Um painel sobreposto
deixaria o participante de óculos sem nível, sem instrução e sem
controles.

% [FIGURA] o console dentro da cena
O painel acompanha o corpo do participante, e não a sua cabeça:
desloca-se com ele quando caminha, mas não gira junto com o olhar.
Interface fixada ao campo de visão é reconhecidamente causa de
desconforto em Realidade Virtual.

\subsubsection{Dois arranjos, determinados pelo campo de visão}

O posicionamento da interface precisou ser calculado, e não estimado. A
câmera do arcabouço mede os seus 80 graus de abertura na
\emph{vertical}; a abertura horizontal resulta da proporção da tela.
Daí decorre uma diferença considerável:

\begin{table}[htbp]
\centering
\caption{Campo de visão horizontal conforme o formato da tela.}
\label{tab:fov}
\begin{tabular}{lrr}
\toprule
\textbf{Tela} & \textbf{Proporção} & \textbf{Abertura horizontal} \\
\midrule
Computador (1280 $\times$ 800) & 1,60 & 106$^\circ$ \\
Celular em pé (390 $\times$ 844) & 0,46 & 42$^\circ$ \\
\bottomrule
\end{tabular}
\end{table}

Com 42 graus de abertura, metade é 21 graus --- e um painel posicionado a
32 graus à esquerda, adequado no computador, fica \emph{inteiramente fora
da tela} em um celular. A aplicação adota, por isso, dois arranjos: no
computador e no óculos o painel fica à esquerda, com os botões abaixo
dele; no celular tudo se centraliza e encolhe, aproveitando a altura que
a tela em pé tem de sobra.

\subsubsection{Dois apontadores}

A seleção de alvos revelou um conflito que merece registro, por ser
próprio de interfaces tridimensionais e não ter equivalente em páginas
convencionais.

Em celulares sem controle de mão, a seleção se dá pelo olhar: a mira
permanece no centro da tela e, após permanecer sobre um alvo por 1,2
segundo, aciona-o. Como no celular o painel de interface é centralizado,
a mira repousava sobre ele e acionava sozinha o botão que estivesse no
centro --- se fosse o de reiniciar, a página recarregava, e recarregava
novamente, em laço.

A correção foi separar os apontadores conforme a sua natureza: a mira
enxerga apenas os estímulos e as portas, enquanto um segundo apontador,
que parte do dedo ou do cursor e nunca do centro da tela, enxerga apenas
a interface. Dentro do óculos, onde não há dedo nem cursor, a mira volta
a enxergar ambos.

Um segundo defeito foi detectado na mesma ocasião: o arcabouço não ignora
objetos invisíveis ao calcular o alcance do raio de seleção, de modo que
um painel oculto permanecia interceptando cliques destinados aos botões
atrás dele. A solução foi distinguir o que \emph{é} interface do que
\emph{pode ser acionado no momento}, atualizando a lista de alvos a cada
abertura ou fechamento de painel.

Ambos os defeitos só se manifestaram quando os testes automatizados
passaram a empregar cliques e toques reais. Até então disparavam o evento
diretamente sobre o elemento, o que contorna o cálculo do raio --- e
teria aprovado os dois casos sem apontar problema. O registro fica como
observação de método: em interface tridimensional, teste por evento
sintético produz falso positivo.

\subsection{Deslocamento}

\begin{table}[htbp]
\centering
\caption{Formas de interação conforme o dispositivo.}
\label{tab:controles}
\small
\begin{tabular}{p{2.2cm}p{2.6cm}p{3.5cm}p{3.6cm}}
\toprule
\textbf{Dispositivo} & \textbf{Olhar} & \textbf{Caminhar} & \textbf{Selecionar} \\
\midrule
Computador & arrastar o cursor & teclas W A S D & clique \\
Celular & girar o aparelho & controle circular na tela & toque \\
Óculos & mover a cabeça & deslocamento físico & raio do controle ou olhar \\
\bottomrule
\end{tabular}
\end{table}

No celular não há teclado, e girar o aparelho permite olhar em volta mas
não mover-se --- o que inviabilizaria tanto a exploração quanto a
descoberta dos exemplares escondidos. Foi implementado, para esse caso,
um controle circular de toque, cujo deslocamento é \emph{relativo ao
olhar} e limitado pelas paredes do cenário, com folga de 45~cm.

A velocidade adotada é de 1,4~m/s, próxima da caminhada humana. Valores
superiores produzem desconforto em Realidade Virtual.

\subsection{Verificação técnica}
\label{sec:verificacao}

% [FIGURA] a ferramenta de verificação em execução
O protótipo conta com ferramenta própria de verificação automatizada,
que percorre as vinte e cinco combinações de animal e nível e confere, em
cada uma: o carregamento do modelo, a escala aplicada, o apoio no piso, o
enquadramento no campo de visão, o carregamento do áudio, a taxa de
quadros sustentada e a ausência de erros de execução.

A ferramenta não é um acessório. Foi ela que detectou o erro de escala da
cobra mencionado anteriormente, e é ela que permite afirmar, com número,
que uma alteração não degradou o que já funcionava.

\subsubsection{Desempenho}

A medição é feita em renderização por software, sem aceleração por placa
de vídeo --- condição deliberadamente desfavorável. Em qualquer aparelho
com placa de vídeo o desempenho tende a atingir o limite da tela.

A Tabela~\ref{tab:desempenho} compara o desempenho antes e depois do
conjunto de otimizações descrito a seguir, na mesma máquina e com a mesma
área renderizada.

\begin{table}[htbp]
\centering
\caption{Quadros por segundo antes e depois das otimizações, com o
número de exemplares entre parênteses.}
\label{tab:desempenho}
\begin{tabular}{llrrrrr}
\toprule
\textbf{Estímulo} & & \textbf{N1} & \textbf{N2} & \textbf{N3} & \textbf{N4} & \textbf{N5} \\
\midrule
Cobra  & antes  & 16 (1) & 11 (3) & 8 (6)  & 6 (10)  & 4 (15) \\
       & depois & 21 (1) & 20 (2) & 19 (3) & 18 (5)  & 16 (8) \\
\addlinespace
Rato   & antes  & 43 (1) & 35 (3) & 26 (6) & 20 (10) & 15 (15) \\
       & depois & 46 (1) & 40 (3) & 31 (6) & 30 (10) & 21 (15) \\
\addlinespace
Barata & antes  & 37 (1) & 35 (3) & 34 (6)  & 32 (10) & 29 (15) \\
       & depois & 37 (1) & 36 (5) & 34 (10) & 32 (16) & 29 (24) \\
\addlinespace
Aranha & antes  & 44 (1) & 37 (3) & 29 (6) & 24 (10) & 19 (15) \\
       & depois & 47 (1) & 41 (3) & 35 (6) & 32 (10) & 25 (15) \\
\addlinespace
Sapo   & antes  & 26 (1) & 22 (3) & 18 (6) & 14 (10) & 11 (15) \\
       & depois & 26 (1) & 24 (3) & 21 (6) & 19 (10) & 14 (15) \\
\bottomrule
\end{tabular}
\end{table}

A média das vinte e cinco combinações passou de 24 para 29 quadros por
segundo, e o número de combinações sem aviso de desempenho subiu de oito
para doze. Em nenhuma das duas medições houve falha.

Três leituras merecem registro:

\begin{itemize}
  \item \textbf{a cozinha passou de 15 para 24 exemplares sem perder
  desempenho} (29 quadros por segundo em ambas as medições, no Nível~5).
  O que paga esse aumento é o comportamento \emph{quieto}, que dispensa
  o reprodutor de animação, somado ao limite de oito exemplares animados
  simultaneamente;

  \item \textbf{a floresta passou de 4 para 16 quadros por segundo no
  Nível~5}, por duas causas somadas: a redução das texturas do modelo de
  cobra de 1024 para 512 pixels e a redução pela metade da quantidade de
  exemplares;

  \item \textbf{a mesma quantidade tornou-se mais barata}: com quinze
  ratos, de 15 para 21 quadros por segundo; com quinze aranhas, de 19
  para 25. Nada mudou no modelo nem no cenário --- a diferença está
  inteiramente no comportamento dos exemplares.
\end{itemize}

\subsubsection{Otimizações aplicadas}

\begin{table}[htbp]
\centering
\caption{Otimizações e respectivo efeito.}
\label{tab:otimizacoes}
\begin{tabular}{p{4.2cm}p{8.4cm}}
\toprule
\textbf{Otimização} & \textbf{Efeito} \\
\midrule
Piso pintado em canvas &
o padrão do piso era composto por 90 peças separadas; passou a ser
pintado uma vez numa textura e repetido: um desenho em lugar de noventa \\
Céu como cor de fundo &
substituição da esfera de céu por cor de fundo da cena \\
Limite de exemplares animados &
no máximo oito reproduzem a animação do arquivo; cada um recalcula a
posição de cada osso a cada quadro \\
Comportamento \emph{quieto} &
exemplares parados não recebem reprodutor de animação \\
Redução de texturas &
cobra: cinco texturas de 1024 para 512 pixels, de 6,13~MB para 3,82~MB \\
Painel em canvas único &
a interface inteira é desenhada numa só textura: um objeto por quadro,
em lugar de nove \\
\bottomrule
\end{tabular}
\end{table}

\subsection{Modelos tridimensionais}

\begin{table}[htbp]
\centering
\caption{Modelos empregados e animações utilizadas.}
\label{tab:modelos}
\begin{tabular}{lrrll}
\toprule
\textbf{Arquivo} & \textbf{Tamanho} & \textbf{Triângulos} & \textbf{Usadas} & \textbf{Descartadas} \\
\midrule
cobra.glb  & 3,82 MB & 4.380 & Animation   & --- \\
aranha.glb & 0,43 MB & 2.712 & Idle, Walk, Jump & Attack, Death \\
sapo.glb   & 0,56 MB & 4.920 & Idle, Jump  & Attack, Death \\
rato.glb   & 0,53 MB & 4.004 & Idle, Walk, Jump & Run, Attack, Death \\
barata.glb & 0,52 MB & 576   & idle, walking & Attack, death, flying \\
\bottomrule
\end{tabular}
\end{table}

A reprodução das animações exigiu um componente próprio: o arcabouço
carrega o modelo e as suas animações, mas não as reproduz. O componente
desenvolvido faz a ponte com o reprodutor do Three.js e localiza o clipe
por correspondência parcial de nome --- um modelo pode nomear a sua
animação de repouso como \texttt{SnakeArmature|Snake\_Idle}, e procurar
por \texttt{Idle} já a encontra. Isso evita depender do nome exato, que
varia de modelo para modelo.

\subsection{Áudio}

Cada animal possui um som próprio, posicionado na cena, de modo que o
volume percebido varia com a distância --- a progressão dos níveis também
se ouve. Há ainda um som de ambiente, não posicionado.

O áudio começa desligado e depende de ativação explícita. Não se trata de
opção de interface: navegadores proíbem a reprodução de áudio antes de
uma interação do usuário com a página, e como o participante chega vindo
de outra página, a interação anterior não conta. A preferência é
preservada entre níveis durante a mesma sessão.

\subsection{Limites conhecidos do artefato}

Registram-se, por honestidade metodológica, as limitações do artefato que
são de natureza técnica:

\begin{itemize}
  \item os estímulos pequenos são apresentados em escala maior que a
  real, pelas razões expostas; o fator está documentado, mas a
  apresentação não é fisicamente fiel;

  \item os sons foram sintetizados e não são gravações dos animais;

  \item não há medida fisiológica. A possibilidade de captar frequência
  cardíaca foi investigada e descartada: exige sensor com perfil padrão
  de transmissão, que os relógios de uso comum não expõem, e
  acrescentaria ao estudo uma natureza clínica que o seu escopo não
  comporta;

  \item o desempenho medido refere-se à renderização por software;
  números obtidos no aparelho efetivamente utilizado nas sessões devem
  ser coletados com a mesma ferramenta.
\end{itemize}

%==================== RESULTADOS E DISCUSSÃO ====================
\section{Resultados e discussão}

% ESCREVER DEPOIS DA COLETA.
%
% Regra que vale para a seção inteira: nenhum número aqui é digitado
% à mão. As figuras e tabelas são GERADAS pelas rotinas em
% data/analytics, a partir dos dados já anonimizados (P01, P02, ...).
% É o que torna a análise reprodutível e o que permite recompilar o
% documento sem acesso a dado identificável.

\subsection{Caracterização da amostra}

% \input{tables/tab_amostra}
% Gerar com: idade (média, desvio, faixa), curso, período, experiência
% prévia com Realidade Virtual.
[ESCREVER]

\subsection{Medo declarado por estímulo}

% Do questionário inicial (Apêndice B). Gráfico de barras com a média
% do medo declarado (0--10) para cada um dos cinco animais.
% Serve de linha de base para tudo o que vem depois.
[ESCREVER]

\subsection{Progressão pelos níveis}

% Quantos participantes alcançaram cada nível, e quanto tempo
% permaneceram em cada um. O tempo por nível é informativo: parada
% longa pode indicar tanto exploração quanto hesitação.
[ESCREVER]

\subsection{Comportamento durante a exposição}

% As variáveis do Quadro~\ref{tab:variaveis}: aproximação,
% afastamento, agitação, atenção ao estímulo, exploração, interação.
[ESCREVER]

\subsection{Medo declarado frente ao comportamento observado}

% A análise mais interessante do trabalho, e a que melhor responde à
% questão de pesquisa: quem declarou medo alto manteve maior
% distância? Chegou a níveis mais baixos? Olhou menos para o
% estímulo?
%
% É o cruzamento entre o Apêndice B e as medidas automáticas -- entre
% o que a pessoa disse e o que ela fez.
[ESCREVER]

\subsection{Percepção dos participantes}

% Do questionário final (Apêndice C): presença, realismo, facilidade
% de uso, clareza das instruções, desconforto físico, sensação de
% controle sobre o próprio ritmo.
[ESCREVER]

\subsection{Desempenho nos aparelhos utilizados}

% Rodar ferramentas/teste-sistema.html em cada aparelho usado nas
% sessões e registrar aqui. Os números da Seção 5 referem-se a
% renderização por software e não substituem esta medição.
[ESCREVER]

\subsection{Limitações}

% Ser explícita aqui conta a favor na banca. Candidatas já
% identificadas:
%
%   - amostra pequena e por conveniência;
%   - ausência de grupo de controle;
%   - participantes não necessariamente fóbicos -- o trabalho mede
%     percepção e interação, não eficácia;
%   - sessão única, sem acompanhamento ao longo do tempo;
%   - ausência de medida fisiológica;
%   - estímulos pequenos apresentados em escala maior que a real;
%   - sons sintetizados, não gravados.
[ESCREVER]

%========================== CONCLUSÃO ==========================
\section{Conclusão}

% ESTRUTURA EM TRÊS CAMADAS -- é o padrão adotado pelo orientador, e
% serve para não prometer mais do que o dado sustenta:
%
%   EM TERMOS DE FATOS          o que foi construído e o que se observou
%   EM TERMOS DE INTERPRETAÇÃO  o que os dados permitem afirmar
%   EM TERMOS DE HIPÓTESE       o que fica como conjectura a confirmar
%
% No caso deste trabalho a separação é especialmente importante: tudo
% o que toca eficácia terapêutica pertence à TERCEIRA camada, nunca
% às duas primeiras.
%
% O que já se pode afirmar na camada de FATOS, independentemente da
% coleta, porque está medido:
%   - o artefato foi construído e funciona nas 25 combinações de
%     animal e nível, sem falha na verificação automatizada;
%   - executa em computador, celular e óculos a partir do mesmo
%     endereço, sem instalação;
%   - sustenta 24 exemplares animados mantendo a taxa de quadros,
%     em renderização por software.

[ESCREVER]

\subsection{Contribuições}

% Separar as duas: a de engenharia e a de pesquisa.
%
% De engenharia, já sustentadas:
%   - protótipo de exposição gradual executável em navegador, com a
%     interface construída dentro da cena tridimensional;
%   - ferramenta de verificação automatizada das 25 combinações;
%   - conjunto de decisões de projeto documentadas e justificadas
%     (escala de apresentação, comportamento por exemplar, separação
%     dos apontadores).
%
% De pesquisa: depende da coleta.
[ESCREVER]

\subsection{Trabalhos futuros}

% Candidatos já levantados ao longo do desenvolvimento:
%   - rastreamento de mãos por visão computacional, usando a câmera
%     do aparelho, para permitir aproximação com a própria mão;
%   - medida fisiológica (frequência cardíaca) com sensor de perfil
%     padrão, como cinta torácica;
%   - estudo com participantes com fobia diagnosticada, mediante
%     aprovação em Comitê de Ética em Pesquisa;
%   - acompanhamento ao longo de várias sessões, em lugar de sessão
%     única;
%   - apresentação em escala real, verificando se a perda de
%     visibilidade compromete ou não o estímulo.
[ESCREVER]


%=================== ELEMENTOS PÓS-TEXTUAIS ===================
%========================== REFERÊNCIAS ==========================
\newpage
\section*{Referências}
\addcontentsline{toc}{section}{REFERÊNCIAS}

\begingroup
\singlespacing
\setlength{\parindent}{0pt}
\setlength{\parskip}{0.7\baselineskip}

% ABNT, ordem alfabética. Modelos de cada tipo:
%
% ARTIGO EM PERIÓDICO
% SOBRENOME, N. Título do artigo. \textit{Nome do Periódico}, v. 10, n. 2,
% p. 100--115, 2024.
%
% TRABALHO EM EVENTO
% SOBRENOME, N. Título. In: NOME DO EVENTO, 15., 2024, Cidade.
% \textit{Anais} [...]. Porto Alegre: SBC, 2024. p. 9--15.
%
% LIVRO
% SOBRENOME, N. \textit{Título do livro}: subtítulo. Cidade: Editora, 2023.
%
% PÁGINA NA INTERNET
% A-FRAME. \textit{A-Frame documentation}. Disponível em:
% https://aframe.io/docs/. Acesso em: 9 nov. 2026.

[INSERIR AS REFERÊNCIAS DOS ARTIGOS SELECIONADOS]

\endgroup

%=========================== APÊNDICES ===========================
\newpage
\section*{Apêndice A --- Termo de Consentimento Livre e Esclarecido}
\addcontentsline{toc}{section}{APÊNDICE A --- TCLE}

% RASCUNHO. Conferir com o orientador se a instituição exige submissão
% ao Comitê de Ética em Pesquisa (CEP) e se há modelo próprio de TCLE.
% Este texto é também o que a tela de cadastro da aplicação apresenta.

\begingroup\singlespacing\small

\noindent\textbf{Título do estudo:} \tcctitulo

\noindent\textbf{Pesquisadora:} \tccautor \quad \textbf{Orientador:} \tccorientador

\vspace{0.4cm}
\noindent\textbf{Do que se trata.} Você está sendo convidado(a) a participar de
uma pesquisa de Trabalho de Conclusão de Curso que avalia um protótipo de
Realidade Virtual. O sistema apresenta animais comumente associados a fobias
--- cobra, rato, barata, aranha e sapo --- em ambientes virtuais, organizados em
níveis crescentes de aproximação.

\vspace{0.3cm}
\noindent\textbf{O que não é.} Este sistema \textbf{não é tratamento
psicológico ou médico}, \textbf{não realiza diagnóstico} e \textbf{não
substitui acompanhamento profissional}. A avaliação restringe-se a usabilidade,
imersão, interação e percepção.

\vspace{0.3cm}
\noindent\textbf{O que você fará.} Preencherá um cadastro e um questionário
curto, usará a aplicação pelo tempo que desejar e responderá a um questionário
final. A sessão dura cerca de [N] minutos.

\vspace{0.3cm}
\noindent\textbf{Dados coletados.} E-mail, sexo e data de nascimento; o grau de
medo que você declara para cada animal; e medidas automáticas de uso, como
distância mantida até o estímulo, movimentação e níveis percorridos.

\vspace{0.3cm}
\noindent\textbf{Como seus dados são tratados.} Seu e-mail serve apenas para
identificar o seu cadastro e \textbf{não acompanha os dados de análise}: nos
resultados você é identificado(a) por um código (P01, P02, \ldots). Nenhum dado
que permita identificá-lo(a) será publicado.

\vspace{0.3cm}
\noindent\textbf{Riscos e desconforto.} O contato com imagens de animais pode
causar desconforto. \textbf{Você pode interromper a qualquer momento}, sem
precisar justificar e sem prejuízo de qualquer natureza. O uso de óculos de
Realidade Virtual pode provocar tontura passageira.

\vspace{0.3cm}
\noindent\textbf{Benefícios.} Não há benefício direto. A participação contribui
para a avaliação de uma ferramenta de pesquisa.

\vspace{0.3cm}
\noindent\textbf{Voluntariedade.} A participação é voluntária. Você pode
desistir a qualquer momento e solicitar a exclusão dos seus dados pelo contato
[E-MAIL DA PESQUISADORA].

\vspace{0.6cm}
\noindent Declaro que li e compreendi as informações acima e concordo em
participar.

\vspace{0.8cm}
\noindent\rule{8cm}{0.4pt}\hfill\rule{4cm}{0.4pt}\\
Assinatura do(a) participante \hfill Data

\endgroup

\newpage
\section*{Apêndice B --- Questionário inicial de perfil}
\addcontentsline{toc}{section}{APÊNDICE B --- Questionário inicial}

% Aplicado no cadastro, dentro da própria aplicação.
% Curto de propósito: a pessoa está curiosa para entrar.

\begin{enumerate}
  \item Sexo: ( ) Feminino \quad ( ) Masculino \quad ( ) Outro \quad ( ) Prefiro não informar
  \item Data de nascimento: \rule{3cm}{0.4pt}
  \item Curso e período: \rule{6cm}{0.4pt}
  \item Já usou óculos de Realidade Virtual? ( ) Sim \quad ( ) Não
  \item De 0 a 10, quanto medo você sente de cada animal?
  \begin{itemize}
    \item Cobra: \rule{1.5cm}{0.4pt} \quad Rato: \rule{1.5cm}{0.4pt} \quad
          Barata: \rule{1.5cm}{0.4pt} \quad Aranha: \rule{1.5cm}{0.4pt} \quad
          Sapo: \rule{1.5cm}{0.4pt}
  \end{itemize}
\end{enumerate}

\newpage
\section*{Apêndice C --- Questionário final de percepção}
\addcontentsline{toc}{section}{APÊNDICE C --- Questionário final}

% Respostas de 0 a 10, salvo indicação em contrário.

\begin{enumerate}
  \item O quanto você se sentiu presente dentro do ambiente virtual?
  \item O quanto o animal pareceu real para você?
  \item O quanto foi fácil aprender a usar o sistema?
  \item O quanto foi fácil se deslocar pelo ambiente?
  \item As instruções estavam claras?
  \item Você sentiu desconforto físico (tontura, enjoo)?
  \item O quanto você se sentiu no controle do próprio ritmo?
  \item Em que nível você sentiu maior desconforto? ( ) 1 ( ) 2 ( ) 3 ( ) 4 ( ) 5 ( ) nenhum
  \item Você interromperia a experiência em algum momento? Por quê?
  \item O que você mudaria no sistema?
\end{enumerate}


\end{document}
